% Example lecture deck. Compile: latexmk -xelatex lecture-example.tex
\documentclass[10pt,aspectratio=169,t]{beamer}
\usetheme{upbminimal}

\course{Application Development for Mobile Devices}
\decklabel{Lecture 5}
\title{State Management, Part 1}
\subtitle{Debugging, Stateless vs Stateful, and where state should live}
\author{Dinu-Ștefan Rusu}
\institute{FILS, Universitatea Politehnica București}
\date{Autumn 2026}

\begin{document}

\titleframe

\objectivesframe{
  \cell[\faBug]{Debug on purpose}{Launch and attach DevTools, read the inspector, set breakpoints and step through Dart.}
  \cell[\faLayerGroup]{The three trees}{Widget, Element and RenderObject, and what setState actually triggers.}
  \cell[\faToggleOn]{Stateless vs Stateful}{When a widget really needs a State object, and when it genuinely does not.}
  \cell[\faIcon{share-alt}]{Sharing state}{Lifting state up, InheritedWidget, and why value equality decides rebuilds.}
}

\divider{Part 1 · Debugging}{Seeing what your app is actually doing}[DevTools, breakpoints, and useful logging]

\begin{frame}[eyebrow=Debugging]{The debugging toolbox}
  \vskip 2mm
  \begin{grid}[3]
    \cell[\faIcon{sync-alt}]{Hot reload}{Injects changed source into the running VM and rebuilds. State objects survive.}
    \cell[\faPowerOff]{Hot restart}{Rebuilds from scratch and drops all state. Use it whenever initState or a global changed.}
    \cell[\faTerminal]{debugPrint \& assert}{Cheap, always available, and compiled out of release builds entirely.}
    \cell[\faPause]{Breakpoints}{Stop on a line, inspect every variable in scope, step through the frame.}
    \cell[\faIcon{tachometer-alt}]{DevTools}{Inspector, performance, CPU, memory, network: a browser tab attached to your app.}
    \cell[\faExclamationCircle]{Read the error}{Flutter's exception boxes name the widget and the constraint that failed.}
  \end{grid}
  \note{Ask the room how many have ever opened DevTools. Usually very few; that is the gap this section closes.}
\end{frame}

\begin{frame}[fragile,eyebrow=Debugging]{print, debugPrint, and assert}
  \begin{columns}[T]
    \begin{column}{0.56\textwidth}
\begin{dart}
print('value: $x');       // logcat truncates
debugPrint('value: $x');  // throttled, lossless
assert(count >= 0, 'count negative');

debugDumpApp();            // widget tree
debugDumpRenderTree();     // sizes + constraints
(*@\hl{debugPrintRebuildDirtyWidgets = true;}@*)
\end{dart}
    \end{column}
    \begin{column}{0.40\textwidth}
      \lead{Why not just print?}{Android's logcat drops long lines. \code{debugPrint} throttles output so nothing disappears.}
      \muted{\code{assert} takes a message. Use it for invariants you want to hear about in the lab.}
    \end{column}
  \end{columns}
  \begin{takeaway}{Logging is a primary debugging tool.}
    It is the only tool that works identically on a simulator, a real phone and a colleague's machine.
  \end{takeaway}
\end{frame}

\begin{frame}[eyebrow=State]{Pick the smallest tool that works}
  \vskip 1mm
  {\usebeamerfont{small}\begin{tabular}{@{}lllp{14mm}p{44mm}@{}}
    \toprule
    \thead{Tool} & \thead{Scope} & \thead{Boilerplate} & \thead{Testable} & \thead{Use when} \\
    \midrule
    setState        & one widget       & none   & partly    & the state never leaves the screen \\
    Lifting up      & a subtree        & low    & yes       & two siblings share one value \\
    InheritedWidget & any descendant   & medium & yes       & a value is read far below its owner \\
    Provider        & any descendant   & low    & yes       & the same, with less ceremony \\
    BLoC / Cubit    & whole feature    & high   & excellent & events and side effects \\
    \bottomrule
  \end{tabular}}
  \vskip 3mm
  \muted{Provider wraps InheritedWidget. Riverpod is its successor, without the BuildContext.}
  \begin{takeaway}{This is Part 1.}
    setState, lifting up and InheritedWidget carry most screens. The last column is Lecture 6.
  \end{takeaway}
\end{frame}

\begin{frame}[eyebrow=Numbers]{Why the frame budget matters}
  \vskip 4mm
  \begin{grid}[3]
    \bignum{16.7 ms}{per frame at 60 Hz}
    \bignum{8.3 ms}{per frame at 120 Hz, most 2026 phones}
    \bignum{1}{missed frame is visible jank}
  \end{grid}
\end{frame}

\begin{frame}[embedded,eyebrow=Embedded,fragile]{Publishing a reading over MQTT}
  \begin{columns}[T]
    \begin{column}{0.58\textwidth}
\begin{golang}
func main() {
    led := machine.LED
    led.Configure(machine.PinConfig{Mode: machine.PinOutput})
    for {
        led.High()
        time.Sleep(200 * time.Millisecond)
        led.Low()
        time.Sleep(800 * time.Millisecond)
    }
}
\end{golang}
    \end{column}
    \begin{column}{0.38\textwidth}
      \lead{Orange marks embedded topics.}{The \code{[embedded]} frame option switches the accent for that slide only.}
      \begin{hint}[Board]
        ESP32-C3, flashed with \code{tinygo flash}.
      \end{hint}
    \end{column}
  \end{columns}
\end{frame}

\begin{frame}[eyebrow=Recap]{Three things to remember}
  \vskip 2mm
  \begin{enumerate}
    \item Widgets are immutable blueprints. \muted{The State object is what persists across builds.}
    \item setState marks one Element dirty. \muted{Only that subtree rebuilds on the next frame.}
    \item Value equality decides rebuilds. \muted{That is the whole reason Equatable exists.}
  \end{enumerate}
  \vskip 5mm
  \kv{Further reading}{\link{https://docs.flutter.dev/ui/interactivity}{docs.flutter.dev/ui/interactivity} · \link{https://docs.flutter.dev/tools/devtools}{DevTools guide}}
\end{frame}

\closingframe{Questions?}{dinu\_stefan.rusu@upb.ro · Microsoft Teams: course channel}

\end{document}
